\documentclass[book.tex]{subfiles}

\begin{document}
\chapter{Case study: Identifying Photovoltaic Materials}

\label{chap:photovoltaic_case}
\noindent\rule{11cm}{0.4pt} \\
\textbf{Prerequisites:} \autoref{chap:graphconvs}, \autoref{chap:equivariant_networks}, \autoref{chap:molecular_dynamics},  \autoref{chap:dft} \\
\textbf{Difficulty Level:} ** \\
\noindent\rule{11cm}{0.4pt}
\vspace{5mm}

TODO: Physika case study

The case study will focus on developing deep learning models to predict properties relevant for identifying promising photovoltaic materials. Chapter 1's section on photovoltaics provides an overview of materials development in the field of photovoltaics and highlights milestones over the past few decades. 
\begin{figure}[h]
    \centering
    \includegraphics[width=0.7\linewidth]{figures/gapParityNew.png}
    \caption{Parity plots for bandgap predictions using different featurization techniques. A: ECFP, B: Graph Convolution, and C: Weave }
    \label{fig:gapParity}
\end{figure}

\subsection{Data}
Our data is sourced from \textbf{Citrination}, the open database and analytics platform of Citrine Informatics. The dataset is sourced from the Harvard Organic Photovoltaic Dataset (HOPV15), which was retrieved through the Citrination API. The dataset is a collation of experimental photovoltaic data through literature and quantum-chemical calculations performed over a range of prospective organic photovoltaic materials. There are over 2.3 million data entries, each with their respective identifying notation, also called as SMILES (Simplified Molecular Input Line Entry System), and material properties, including HOMO and LUMO levels, HOMO-LUMO gap, power conversion efficiency (PCE), Open Circuit Voltage (OCV) and short-circuit current density (SCCD). In Figure \ref{fig:properties}, the distribution of our dataset is shown for each property that was extracted. We can infer that most of them follow a normal distribution and thus be easily captured by an accurate ML model.  
\begin{figure}[h]
    \centering
    \includegraphics[width=0.7\linewidth]{figures/homoParityNew.png}
    \caption{Parity plots for HOMO level predictions using different featurization techniques. A: ECFP, B: Graph Convolution, and C: Weave }
    \label{fig:gapParity}
\end{figure}
% \begin{figure}[ht]
%     \centering
%     \includegraphics[width=\linewidth, height=15cm]{figures/raw_data_properties_histograms.png}
%     \caption{Distribution of extracted properties}
%     \label{fig:properties}
% \end{figure}

\subsection{Methods and Model Optimization}
Similar to the previous case study, MoleculeNet\cite{Pande} in conjunction with the open-source package DeepChem is employed for featurization and model development for predictions of relevant properties in the context of photovoltaic materials. To bridge the gap between material characteristics and input needed for Deepchem, RDKit, a chemical informatics library for Python, is used to generate descriptors useful for machine learning from molecules as input. 

The study can be broken down into three parts: first, accessing and locally visualizing the dataset from Citrination; using DeepChem and RDKit tools to perform the required machine learning techniques for a smaller part of the dataset; and lastly, using the full dataset with superior computational resources (Google Cloud) to create new Machine Learning (ML) models. 

From the DeepChem library, we used three different featurization methods to train ML models on four properties that are most relevant in the context of photovoltaics: PCE, OCV, SCCD and gap. The featurization methods chosen are similar to the previous case study: Extended-Connectivity Fingerprints (ECFP), Graph Convolutions (GC), and Weave. \cite{Wu2018MoleculeNet:Learning} The whole data set was split into a 80:10:10 ratio for training, validation, and testing subsets. Finally, a brief analysis of loss per epoch for the GC and Weave techniques in SCCD training is presented as a way to study the particularly bad performance of these methods in predicting some of the properties of interest. 

The ECFP featurization technique has been implemented alongside a Random Forest regressor to train ML models that predict the properties of interest, one model per property.
The GC featurization method was implemented alongside a neural network with 3 hidden convolutional layers and two fully connected layers in a regression scheme to accurately predict the relevant property. Max pooling, zero padding, batch normalization and Relu activation function is applied in every hidden convolutional layer, while tanh activation function and batch normalization is applied in the fully connected (dense) layers. For the training process, the learning rate was selected to be 0.001, batch size of 250 and number of epochs were 10 for all 2.4 million data points available from Citrination. Finally, Adam optimizer was selected to minimize the Mean Squared Error (MSE) loss function after applying a normalization transformation to all available data. This featurization method was implemented alongside a neural network, the structure and hyper parameters of which is identical with the one utilized in Graph Convolution model explained above.


% \textbf{Extended-Connectivity Fingerprints (ECFP)}
% The featurization decomposes a molecule into submodules originated from heavy atoms, each assigned a unique identifier. These segments are extend through bonds to form larger substructures and corresponding identifiers. After hashing all these substructures into a fixed length
% binary fingerprint, the representation contains information about topological characteristics of the molecule, which enables it to be applied to tasks such as similarity searching and activity prediction. \\

% The ECFP featurization technique has been implemented alongside a Random Forest regressor to train ML models that predict the properties of interest, one model per property.

% \textbf{Graph Convolutions (GC)}
% The GC featurization computes an initial feature vector and a neighbor list for each atom. The feature vector summarizes the atom's local chemical environment, including atom-types, hybridization types, and valence structures. Neighbor lists represent connectivity of the whole molecule, which are further processed in each model to generate graph structures. \\

% This featurization method was implemented alongside a neural network with 3 hidden convolutional layers and two fully connected layers in a regression scheme to accurately predict the relevant property. Max pooling, zero padding, batch normalization and Relu activation function is applied in every hidden convolutional layer, while tanh activation function and batch normalization is applied in the fully connected (dense) layers. For the training process, the learning rate was selected to be 0.001, batch size of 250 and number of epochs were 10 for all 2.4 million data points available from Citrination. Finally, Adam optimizer was selected to minimize the Mean Squared Error (MSE) loss function after applying a normalization transformation to all available data. 

% \textbf{Weave}
% Similar to GC, the Weave featurization encodes both local chemical environment and connectivity of atoms in a molecule. Atomic feature vectors are exactly the same, while connectivity uses more detailed pair
% features instead of neighbor listing. The Weave featurization calculates a feature vector for each pair of atoms in the molecule, including bond properties (if directly connected), graph distance and ring info, forming a feature matrix. The method supports graph-based models that utilize properties of both nodes (atoms) and edges (bonds). This featurization method was implemented alongside a neural network, the structure and hyper parameters of which is identical with the one utilized in Graph Convolution model explained above.
Multitask training and prediction was used for Weave model and GraphConv model. 

Hyperparameter Optimization (Model Learning): Grid search method was used for hyperparameter optimization. Number of training epochs ($n_epochs$), batch size and learning rate were the parameters over which the optimization was performed.

Weave model configuration: $N_epochs = 150$, $batch size = 16$, and learning rate $= 0.0001$ were identified as the optimal set of hyperparameters for the Weave model. DeepChem’s default Weave model architecture was used: 2 Weave layer + dense layer + batchnorm layer + WeaveGather layer + output dense layer.


\begin{table}[!htbp]
\centering
\caption{Performance with the Weave Model.}
\begin{tabular}{*6c}
\toprule
Evaluation Set  & Size &  \multicolumn{2}{c}{Gap (eV)} & \multicolumn{2}{c}{HOMO Level (eV)}\\
\midrule
     &    & RMSE    & MAE   & RMSE & MAE\\
\midrule
Training&1840823    &  0.035 & 0.027    & 0.025
   &  0.018\\
Validation& 230103   &  0.035 & 0.027   & 0.025  & 0.019\\
Testing& 230103   &  0.035 & 0.027   &  0.025  & 0.018\\
\bottomrule
\end{tabular}
\end{table}


GraphConv Model Configuration: $N_epochs$ = 100, batch size = 8, and learning rate = $0.001$ were identified as the optimal set of hyperparameters for the GraphConv model. DeepChem’s default GraphConv model architecture was used:  2 graph convolutional layers + dense layer + batchnorm layer + graphGather layer + output dense layer.

\begin{table}[!htbp]
\centering
\caption{Performance the GraphConv Model.}
\begin{tabular}{*6c}
\toprule
Evaluation Set  & Size &  \multicolumn{2}{c}{Gap (eV)} & \multicolumn{2}{c}{HOMO Level (eV)}\\
\midrule
     &    & RMSE    & MAE   & RMSE & MAE\\
\midrule
Training&1840823    &  0.034 & 0.026    & 0.023
   &  0.017\\
Validation& 230103   &  0.034 & 0.026   & 0.023  & 0.017\\
Testing& 230103   &  0.034 & 0.026   &  0.023  & 0.017\\
\bottomrule
\end{tabular}
\end{table}


Homo task and Gap task were conducted by two separate ECFP models (ECFP is not supported with multitasking). 

ECFP Model Configuration: number of estimators = 250 was selected as the hyperparameter for Homo task. Everything else followed DeepChem’s default model setting.

\begin{table}[!htbp]
\centering
\caption{Performance the GraphConv Model.}
\begin{tabular}{*6c}
\toprule
Evaluation Set  & Size &  \multicolumn{2}{c}{Gap (eV)} & \multicolumn{2}{c}{HOMO Level (eV)}\\
\midrule
     &    & RMSE    & MAE   & RMSE & MAE\\
\midrule
Training&1840823    &  0.020 & 0.013    & 0.017
   &  0.011\\
Validation& 230103   &  0.049 & 0.032   & 0.042  & 0.028\\
Testing& 230103   &  0.049 & 0.032   &  0.042  & 0.028\\
\bottomrule
\end{tabular}
\end{table}

\subsection{Results}
Each of the three featurization methods were used to train ML models to predict the four properties of interest, yielding a total of twelve different models. These were evaluated over testing data sets that correspond to 10\% of our input data. A summary of the prediction error distributions of each model is shown in Table \ref{tab:results}.  



% \begin{table}[ht]
%     \centering
%     \includegraphics[width=\linewidth, height=16cm]{figures/table.png}
%     \caption{Histograms of the prediction error distribution for each on the three different featurization methods}
%     \label{tab:results}
% \end{table}


% \begin{table}[ht]
%     \centering
%     \includegraphics[width=\linewidth, height=16cm]{figures/gapParity.png}
%     \caption{Parity plots for predictions using different featurization techniques. A: ECFP, B: Graph Convolution, and C: Weave }
%     \label{tab:results}
% \end{table}

Table \ref{tab:results} shows that, in general, the models based on the ECFP featurization technique outperform those reliant on GC and Weave. This result is a little surprising, since the GC and Weave methods are more involved, both conceptually and computationally. It would be expected that the more complex a model is, the better the results. Furthermore, the errors associated with the ECFP models are very small compared to their respective data sets, indicating that this featurization technique yields accurate results. 

\subsection{Discussion}
The SCCD predictions using GC and Weave techniques had their losses (in non-normalized SSE) measured as a function of epoch used for training, as one can see in  Figure \ref{fig:lossEpoch}. As it can be observed, the GC model seems to have reached reasonable convergence in the losses per epoch plot, while the Weave model apparently either needs a higher number of epochs for a more accurate prediction estimation, or it is not a good model for predicting this property with the architecture used (as will be further elaborated in the following paragraph). Additionally, the models trained on the GC and Weave featurizations can have their performance enhanced by adding more layers to the neural networks and by manually tuning the hyper-parameters.
% \begin{figure}[ht]
%     \centering
%     \includegraphics[width=\linewidth, height=9cm]{figures/learning_per_epoch_SCCD.png}
%     \caption{Loss per epoch plot for the GC and Weave based models to predict SCCD}
%     \label{fig:lossEpoch}
% \end{figure}

\subsection{Conclusions}
As debated above, among the three different featurization methods, those models based on the ECFP  technique are the cheapest, from a computational standpoint, and have higher accuracy. Approaches to enhance the performance of GC and weave based models are presented albeit with a higher computational cost. An in-depth analysis of possible correlations present in the data (both in terms of property-property correlations as well as molecular structure-property correlations) should be conducted in order to add intuition to models we have been using as a black box. Furthermore, multitask models to predict all of the properties of interest, using the same featurization techniques as the ones in this report, should be trained to capture these possible correlations. For example, supppose PCE and OCV are related. Models based on GC are quite accurate at predicting OCV, but perform poorly at predicting PCE. However, if there was only one GC-based model trained to predict both of them together, it provides an opportunity to "learn" the correlation, and have a better performance on its PCE predictions. 


\end{document}